\chapter{Quantizations }
\label{qnt}

\section{Deformation Quantization}

Deformation quantization replaces the pointwise product of functions on phase space by a star product
\begin{equation}
f \star g = fg + \frac{i\hbar}{2}{f,g} + O(\hbar^2).
\end{equation}
Quantum mechanics is then viewed as a deformation of the classical observable algebra rather than as an operator theory ab initio.

\section{Geometric Quantization}

Geometric quantization begins with a complex line bundle $L \to M$ equipped with a connection $\nabla$ satisfying
\begin{equation}
\mathrm{curv}(\nabla) = \frac{1}{\hbar}\omega.
\end{equation}
This requires the integrality condition $[\omega]/2\pi\hbar \in H^2(M,\mathbb{Z})$.

To obtain a physical Hilbert space, one introduces a polarization selecting a subset of compatible observables. The choice of polarization is generally non-unique and may not exist globally.

\section{Geometric Quantization}

Deformation quantization modifies the algebra. 
Geometric quantization instead attempts to construct the Hilbert space 
directly from symplectic geometry.

\subsection{Prequantization}

Given $(M,\omega)$ with integral symplectic form:

\[
[\omega] \in H^2(M,\mathbb Z),
\]

there exists a Hermitian line bundle $L \to M$ with curvature:

\[
F_\nabla = -i\omega.
\]

Prequantum Hilbert space:

\[
\mathcal H_{pre} = L^2(M,L).
\]

Prequantum operator:

\[
\hat f = -i\hbar \nabla_{X_f} + f.
\]

However:

\[
[\hat f, \hat g] = i\hbar \widehat{\{f,g\}}.
\]

This space is too large.

\subsection{Polarization}

To reduce degrees of freedom, choose a polarization 
(e.g., vertical polarization in $T^*Q$). 

Wavefunctions become sections constant along half of phase space directions.

Example:
\[
\mathcal H = L^2(Q).
\]

\subsection{Metaplectic Correction}

For consistency of half-forms, one introduces the metaplectic correction.

\subsection{Non-Functoriality}

Quantization is not a functor:

\[
\text{Symplectic manifolds}
\not\longrightarrow
\text{Hilbert spaces}
\]

in a structure-preserving categorical sense.

Obstructions include:

\begin{itemize}
\item Groenewold–Van Hove no-go theorem
\item Dependence on polarization
\item Global topological constraints
\end{itemize}

Thus quantization is not canonical.

%%%%%%%%%%%%%%%%%%%%%%%%%%%%%%%%%%%%%%%%%%%%%%%%%%%%%%%%%%%%
\section{Path Integral Quantization}

In the path-integral approach, transition amplitudes are formally expressed as
\begin{equation}
\langle q_f,t_f | q_i,t_i \rangle = \int \mathcal{D}q, e^{\frac{i}{\hbar}S[q]},
\end{equation}
where $S$ is the classical action. While powerful and intuitive, this formulation typically requires auxiliary constructions for mathematical rigor.

\section{Stochastic quantization}
\label{stochastic-quantization}

\section{Tomographic Quantization}
\label{tomographic-quantization}

